\chapter{Data Architecture and Database System}
\label{ch:data-architecture}

The Wolverine intelligence platform relies on a heterogeneous database strategy designed to isolate adversarial payload collection, schema normalization, and cryptographic auditability. This chapter specifies the data lifecycle from raw extraction to relational persistence, detailing the Prisma schema, source platform topologies, and the exactly-once event processing patterns that guarantee data consistency.

\section{Data Model Overview}
\label{sec:data-overview}
The logical data model transitions strictly from raw, schema-less states to normalized, graph-projected forms. While earlier iterations of this manuscript and some theoretical designs refer to the canonical data as the ``CanonicalRecord'', the implementation relies on the \texttt{NormalizedRecord} Prisma model to serve this exact purpose. The primary entities map to actual Prisma models defined in \texttt{wolverine/prisma/schema.prisma}, representing a robust domain-driven design.

\subsection{Prisma Schema Entities}
\label{sec:prisma-schema}

The database schema is structured into three primary domains: operational ingestion, canonical analysis, and graph projection. Each entity serves a distinct role within this architecture.

\begin{table}[h]
\centering
\resizebox{\textwidth}{!}{
\begin{tabular}{p{0.20\textwidth} p{0.15\textwidth} p{0.65\textwidth}}
\toprule
\textbf{Entity} & \textbf{Type} & \textbf{Purpose} \\
\midrule
\texttt{Source} & Dimension & Represents target sites (Atlas, Briar, Cinder, Drift, Ember). \\
\texttt{CollectionRun} & Operational & Tracks data ingestion runs and batch metadata. \\
\texttt{Capture} & Immutable Fact & Immutable captured payloads (\texttt{locator}, \texttt{contentHash}, \texttt{objectKey}). \\
\texttt{NormalizedRecord} & Core Domain & Parsed canonical data (\texttt{recordData} JSON). Referred to conceptually as CanonicalRecord. \\
\texttt{Quarantine} & Operational & Failed parses stored with their \texttt{errorCode} for analyst review. \\
\texttt{ResolutionCandidate} & Core Domain & Resolved alias bridges containing exact \texttt{feature\_scores} JSON and \texttt{state} (\texttt{pending}, \texttt{linked}, \texttt{review}, \texttt{rejected}). \\
\texttt{GraphNode} & Projection & Persistent graph projection nodes with temporal bounds (\texttt{firstSeenAt}, \texttt{lastSeenAt}). \\
\texttt{GraphEdge} & Projection & Persistent graph edges materialized from BFS traversals. \\
\texttt{Outbox} & Infrastructure & Asynchronous event dispatching for exactly-once processing. \\
\texttt{ProcessedMessage} & Infrastructure & Tracks which events have been consumed to enforce idempotency. \\
\bottomrule
\end{tabular}
}
\caption{Prisma Schema Entities and Domain Roles}
\label{tab:prisma-entities}
\end{table}

\paragraph{Source and Collection Metadata}
The \texttt{Source} entity represents the origin of the intelligence, mapping directly to the distinct adversary platforms being monitored. \texttt{CollectionRun} entities organize individual execution iterations of the scraper infrastructure, aggregating metadata such as timestamps, success rates, and volume statistics to measure pipeline health.

\paragraph{Immutable Payloads}
The \texttt{Capture} model is arguably the most critical operational component, serving as the immutable fact table for raw ingestion. Every HTTP response payload extracted during a \texttt{CollectionRun} is stored durably in MinIO. The \texttt{Capture} model records its \texttt{objectKey} (location in MinIO), the \texttt{locator} (the URL or API endpoint), and a \texttt{contentHash} (SHA-256) to ensure the payload is completely tamper-evident.

\paragraph{Domain Entities}
The \texttt{NormalizedRecord} (conceptually the Canonical Record) stores the standardized interpretation of the raw data inside its \texttt{recordData} JSON column. If normalization fails, the raw data bypasses this entity and is shunted into \texttt{Quarantine}, stamped with an \texttt{errorCode} for human analyst review. \texttt{ResolutionCandidate} bridges these records, providing the exact \texttt{feature\_scores} JSON output generated by the heuristic engine and state management (\texttt{pending}, \texttt{linked}, \texttt{review}, \texttt{rejected}).

\paragraph{Graph and Outbox Infrastructure}
\texttt{GraphNode} and \texttt{GraphEdge} represent the intermediate projected state of the domain entities, optimized for bulk loading into Neo4j. The \texttt{Outbox} and \texttt{ProcessedMessage} tables are strictly infrastructure-level elements designed to guarantee distributed consistency, which is elaborated upon in Section \ref{sec:data-consistency}.

\section{Heterogeneous Source Schemas}
\label{sec:data-source-schemas}
The system ingests data from five distinct, simulated adversarial platforms, each implementing deliberate schema variations to challenge the normalization pipeline.

\begin{table}[h]
\centering
\begin{tabular}{p{0.15\textwidth} p{0.15\textwidth} p{0.15\textwidth} p{0.45\textwidth}}
\toprule
\textbf{Platform} & \textbf{Engine} & \textbf{Format} & \textbf{Schema Characteristics} \\
\midrule
\textbf{Atlas Market} & PostgreSQL & JSON & standard relational payloads; \texttt{vendor\_id} mapped directly. \\
\textbf{Briar Bazaar} & PostgreSQL & HTML & Unstructured DOM. Parsers extract via CSS selectors (e.g., \texttt{\#vendor-handle}). \\
\textbf{Cinder Exchange} & MySQL 8.4 & JSON:API & Nested JSON objects (\texttt{data.attributes}). Identifiers are UUIDs. \\
\textbf{Drift Forum} & PostgreSQL & JSON & Flat structures emphasizing timestamp variations (Unix epochs). \\
\textbf{Ember Commons} & PostgreSQL & GraphQL & Deeply nested \texttt{data.feed.edges[].node} traversal required. \\
\bottomrule
\end{tabular}
\caption{Adversarial Source Schema Variations}
\label{tab:source-schemas}
\end{table}

\section{Canonicalization Pipeline}
\label{sec:data-pipeline}
The transformation from source to canonical proceeds strictly:
\texttt{Raw Capture $\rightarrow$ Schema Adapter $\rightarrow$ Extraction $\rightarrow$ Normalization $\rightarrow$ Hashing $\rightarrow$ Persistence}.

The system utilizes \texttt{processor.ts} to enforce rigorous standardization across all \texttt{NormalizedRecord} entries. A critical aspect of this pipeline is the cryptographic binding of the record. The SHA-256 content hashing is performed via \texttt{cryptoHash()} during processing. This ensures that every canonicalized \texttt{NormalizedRecord} maintains an unbreakable cryptographic link to its raw \texttt{Capture}. If a parser fails, the payload is shunted into the \texttt{Quarantine} table with a specific \texttt{errorCode}, preventing schema corruption in the primary analytical tables.

\textit{Illustrative synthetic example — not an empirical observation:}
\begin{enumerate}
    \item \textbf{Input:} A Cinder Exchange JSON:API payload containing \texttt{"createdAt": 1737331200} (Unix epoch) is saved as a \texttt{Capture}.
    \item \textbf{Extraction:} The parser extracts the integer and other relevant attributes.
    \item \textbf{Normalization:} The parser coerces the integer to an ISO-8601 string, mapping the structure into the unified \texttt{recordData} JSON format for the \texttt{NormalizedRecord}.
    \item \textbf{Hashing:} \texttt{cryptoHash()} digests the original JSON byte array, generating the \texttt{contentHash} (e.g., \texttt{e3b0c442...}).
    \item \textbf{Output:} The resulting \texttt{NormalizedRecord} and its associated \texttt{Outbox} event are persisted atomically to PostgreSQL.
\end{enumerate}

\section{Provenance and Integrity}
\label{sec:data-provenance}
The \texttt{contentHash} guarantees cryptographic provenance. While the \texttt{NormalizedRecord} is heavily mutated (normalized) for analysis, an auditor can always retrieve the exact historical byte representation from \texttt{minio} using the hash (via \texttt{objectKey}), proving the intelligence finding was derived from genuine collection rather than pipeline hallucination. This provides integrity, though in a synthetic environment, it does not guarantee real-world authenticity.

\section{Database and Storage Topology}
\label{sec:data-storage-topology}
The persistence layer spans multiple engines to optimize specific operational workloads. The infrastructure is deployed via Docker Compose, exposing specific ports to facilitate communication and debugging.

\begin{table}[h]
\centering
\resizebox{\textwidth}{!}{
\begin{tabular}{p{0.20\textwidth} p{0.15\textwidth} p{0.15\textwidth} p{0.30\textwidth} p{0.15\textwidth}}
\toprule
\textbf{Service} & \textbf{Engine} & \textbf{Ports} & \textbf{Stored Data} & \textbf{Trust Level} \\
\midrule
\texttt{postgres-sites} & PostgreSQL & 5432 (int) & Atlas, Briar, Drift, Ember source DBs. & Untrusted \\
\texttt{mysql-site-c} & MySQL & 3306 (int) & Cinder source DB. & Untrusted \\
\texttt{postgres-wdb} & PostgreSQL & 5432:5433 & Normalized Records, Linkages. & Semi-Trusted \\
\texttt{postgres-truth} & PostgreSQL & 5432:5434 & Ground Truth assignments. & Strictly Trusted \\
\texttt{minio} & Object & 9000, 9001 & Immutable raw HTTP payload captures. & Semi-Trusted \\
\texttt{redis} & In-Memory & 6379 & Pub/Sub event streams. & Semi-Trusted \\
\texttt{neo4j} & Graph & 7474, 7687 & Projected BFS adjacency matrices. & Semi-Trusted \\
\bottomrule
\end{tabular}
}
\caption{Database Inventory, Topology, and Port Mappings}
\label{tab:database-inventory}
\end{table}

\section{Data Consistency and the Outbox Pattern}
\label{sec:data-consistency}
Maintaining distributed consistency between the primary relational datastore (\texttt{postgres-wdb}) and downstream analytical systems (\texttt{neo4j} graph projections, notification services) requires careful orchestration. The Wolverine architecture implements the transactional outbox pattern to achieve exactly-once event processing guarantees \cite{kleppmann2017designing}.

When a business transaction occurs—such as the creation of a \texttt{ResolutionCandidate} or a new \texttt{NormalizedRecord}—the system does not immediately call external services or directly publish to the message broker. Instead, it leverages the Prisma client to atomically insert the domain entity and a corresponding event record into the \texttt{Outbox} table within the same database transaction.

The flow operates as follows:
\begin{enumerate}
    \item \textbf{Atomic Write:} The domain entity (e.g., \texttt{NormalizedRecord}) and an \texttt{Outbox} event are committed simultaneously. If the transaction fails, neither is saved.
    \item \textbf{Signaling:} A Redis Pub/Sub notification is triggered to wake up the asynchronous relay worker.
    \item \textbf{Relay Processing:} The worker polls the \texttt{Outbox} table, ensuring it only reads from the definitive source of truth.
    \item \textbf{Idempotency:} Before processing an event, the worker checks the \texttt{ProcessedMessage} table. If the event ID exists, it is skipped. If not, the worker executes the downstream logic (e.g., updating \texttt{GraphNode} and \texttt{GraphEdge} models) and records the event ID in \texttt{ProcessedMessage} in a single transaction.
\end{enumerate}

This architecture completely decouples the synchronous web request lifecycle from heavy graph projection workloads, guaranteeing that the \texttt{neo4j} graph visualization is eventually consistent without risking message loss during transient network partitions. The outbox processor continuously tails the table and leverages row-level locking to allow concurrent execution of multiple relay workers without duplicate processing.

\section{Data Lifecycle}
\label{sec:data-lifecycle}
The data lifecycle progresses through distinct stages, visualizing the journey from adversarial platforms to the analyst dashboard. This progression is governed by strict boundaries:
\begin{itemize}
    \item \textbf{Ingestion:} The \texttt{CollectionRun} coordinates distributed scrapers fetching payloads from \texttt{Source} platforms, dropping them into \texttt{minio} and creating a \texttt{Capture}.
    \item \textbf{Processing:} Asynchronous workers parse the \texttt{Capture}, producing either a \texttt{NormalizedRecord} or a \texttt{Quarantine} entry.
    \item \textbf{Resolution:} The entity resolver compares \texttt{NormalizedRecord} entries, generating a \texttt{ResolutionCandidate} with detailed \texttt{feature\_scores} when heuristic thresholds are met.
    \item \textbf{Projection:} Approved \texttt{ResolutionCandidate} linkages trigger outbox events that materialize into \texttt{GraphNode} and \texttt{GraphEdge} records, ultimately synced to Neo4j.
\end{itemize}
These sequential transformations ensure that all high-level analytical insights are directly traceable back to raw, immutable collection events.

\section{WolverineDB Integration}
\label{sec:data-wolverinedb}
The cryptographic integrity layer, originally scoped within this general data architecture, has been decoupled into a dedicated distributed ledger middleware known as WolverineDB. The exact specifications, Merkle tree implementations, replay protection mechanisms, and state invariants of WolverineDB are detailed comprehensively in \Cref{ch:wolverinedb}.

\section{Database Failure Modes}
\label{sec:data-failures}
\begin{itemize}
    \item \textbf{Neo4j Unavailable:} \texttt{analyst-ui} graph queries fail gracefully. Linkage calculation continues unabated in PostgreSQL because \texttt{GraphNode} and \texttt{GraphEdge} projections act as resilient intermediaries.
    \item \textbf{MinIO Unreachable:} Normalization aborts completely. Without a verifiable payload hash, no canonical record is permitted to exist.
    \item \textbf{Malformed Payloads:} The database adapter rejects constraints, throwing a schema validation exception back to the API and routing the payload to the \texttt{Quarantine} table.
\end{itemize}

\section{Data Architecture Trade-offs}
\label{sec:data-tradeoffs}
The primary trade-off in the data architecture is prototype simplicity versus distributed scalability. Using \texttt{neo4j} merely as a dump target for in-memory BFS projections (rather than executing native Cypher queries) artificially bounds the network scale ($N \le 500$), but guarantees determinism and rapid execution for demonstration purposes. Furthermore, the reliance on a completely isolated synthetic Truth Vault prevents the system from encountering the chaotic, unclassifiable noise inherent to real-world threat telemetry.
